\newpage

\section{Number Theory}
\label{sec:numbertheory}

\subsection{Multiples and divisors}
\label{sec:7-1}
% CN 7.1

\begin{definition}[Divisibility]
For $d\in\Z\setminus\{0\}$ and $n\in\Z$, $d$ \emph{divides} $n$, written $d\mid n$, if $n=qd$ for some $q\in\Z$. Then $d$ is a \emph{divisor} (or \emph{factor}) of $n$, and $n$ is a \emph{multiple} of $d$; otherwise we write $d\nmid n$. The set of all multiples of $d$ is
\[
d\Z=\{kd : k\in\Z\}.
\]
A divisor of $n$ is \emph{proper} if it is neither $1$ nor $n$.
\end{definition}

\begin{example}
$2\Z$ is the set of even integers, and $10\Z=\{\dots,-20,-10,0,10,20,\dots\}$.
\end{example}

\begin{note}[Zero is a special case]
Every $d\ne0$ divides $0$, since $0=0\cdot d$. So whenever divisors are involved, $0$ often needs separate treatment.
\end{note}

\begin{proposition}[Linearity of divisibility]
\label{thm:divlinear}
For all $d,m,n,x,y\in\Z$ with $d\ne0$,
\[
d\mid m \ \wedge\ d\mid n \quad\Longrightarrow\quad d\mid (xm+yn).
\]
In particular $d\mid(m+n)$ and $d\mid(m-n)$.
\end{proposition}
\begin{proof}
Write $m=q_1d$ and $n=q_2d$. Then $xm+yn=(xq_1+yq_2)d$, and $xq_1+yq_2\in\Z$, so $d\mid(xm+yn)$.
\end{proof}

\begin{note}[The converse fails]
$2\mid(3+5)$, but $2\nmid3$ and $2\nmid5$. Knowing that $d$ divides a sum tells you nothing about the individual terms.
\end{note}

\begin{note}[Divisibility tests from the last digits]
Some divisors can be detected from the decimal representation alone:
\begin{itemize}
    \item $10\mid n$ iff $n$ ends in $0$;
    \item $5\mid n$ iff $n$ ends in $0$ or $5$;
    \item $2\mid n$ iff $n$ ends in one of $0,2,4,6,8$;
    \item $4\mid n$ iff the last two digits of $n$ form a multiple of $4$.
\end{itemize}
Only a fixed number of digits at the end matter, however large $n$ is (Exercise 7.9).
\end{note}

\subsection{Prime numbers}
\label{sec:7-2}
% CN 7.2

\begin{definition}[Prime, composite]
An integer $p>1$ is \emph{prime} if it has no proper divisors. An integer $n>1$ that is not prime is \emph{composite}.
\end{definition}

\begin{note}[$1$ is neither]
$1$ is neither prime nor composite. Excluding it is what makes prime factorisations unique (Theorem~\ref{thm:fta}): otherwise $6=2\cdot3=1\cdot2\cdot3=1\cdot1\cdot2\cdot3=\cdots$.
\end{note}

\begin{theorem}[Existence of prime factorisation]
\label{thm:primefactexists}
Every integer $n\ge2$ is a product of primes.
\end{theorem}
\begin{proof}
We use strong induction on $n$.

\textbf{Base case} ($n=2$): $2$ is prime, so it is a product of one prime.

\textbf{Inductive step}: Let $n\ge2$ and assume every $k$ with $2\le k\le n$ is a product of primes. If $n+1$ is prime, it is a product of one prime. Otherwise $n+1$ has a proper divisor $a$, so $n+1=ab$ with $1<a,b<n+1$ (if $b$ were $1$ or $n+1$, then $a$ would be $n+1$ or $1$). Since $a,b\le n$, the inductive hypothesis makes each of them a product of primes, and hence so is $n+1=ab$.

By the Principle of Strong Induction, the claim holds for all $n\ge2$.
\end{proof}

\bigskip
\begin{theorem}[Infinitude of primes, Euclid]
\label{thm:infprimes}
There are infinitely many primes.
\end{theorem}
\begin{proof}
Suppose, for contradiction, that the primes are exactly $p_1,\dots,p_k$, and let $N=p_1p_2\cdots p_k+1$. By Theorem~\ref{thm:primefactexists}, $N$ has a prime divisor, which must be some $p_i$. But $p_i$ also divides $p_1\cdots p_k$, so by Proposition~\ref{thm:divlinear} it divides $N-p_1\cdots p_k=1$, which is impossible for a prime. So no finite list contains all the primes.
\end{proof}

\begin{proposition}[Trial division bound]
\label{thm:sqrtbound}
If $n>1$ is composite, then $n$ has a prime divisor $p\le\sqrt n$.
\end{proposition}
\begin{proof}
Write $n=ab$ with $1<a\le b<n$. Then $a^2\le ab=n$, so $a\le\sqrt n$. By Theorem~\ref{thm:primefactexists}, $a$ has a prime divisor $p$, and $p\le a\le\sqrt n$ with $p\mid a\mid n$.
\end{proof}

\begin{note}[Using the bound]
To test whether $n$ is prime, it is enough to try dividing by the primes up to $\sqrt n$. For example, $97$ is prime because none of $2,3,5,7$ divides it, and $11^2=121>97$.
\end{note}

\begin{note}[Density of primes (extra)]
The Prime Number Theorem says that $\pi(x)$, the number of primes $\le x$, satisfies $\pi(x)\sim x/\ln x$. So primes thin out, but only slowly: an integer near $x$ is prime with probability about $1/\ln x$. This is why searching randomly for large primes (as needed in \S\ref{sec:7-18}) is practical.
\end{note}

\subsection{Remainders and the mod operation}
\label{sec:7-3}
% CN 7.3

\begin{theorem}[Division algorithm]
\label{thm:division}
For all $n\in\Z$ and $d\in\N$ there exist \emph{unique} $q,r\in\Z$ with
\[
n=qd+r, \qquad 0\le r\le d-1.
\]
\end{theorem}
\begin{proof}
\textbf{Existence}: The set $S=\{n-kd : k\in\Z,\ n-kd\ge0\}$ is nonempty (take $k=-|n|$), so it has a least element $r=n-qd\ge0$. If $r\ge d$, then $r-d=n-(q+1)d$ would be a smaller element of $S$, a contradiction. So $0\le r\le d-1$.

\textbf{Uniqueness}: Suppose $n=q_1d+r_1=q_2d+r_2$ with $0\le r_1,r_2\le d-1$. Then $r_1-r_2=(q_2-q_1)d$, so $d\mid(r_1-r_2)$. But $|r_1-r_2|\le d-1<d$, so $r_1-r_2=0$, and then $q_1=q_2$.
\end{proof}

\begin{definition}[The mod operation]
The unique $r$ of Theorem~\ref{thm:division} is written $n\bmod d$ (read ``$n$ mod $d$''), and the quotient is $q=\lfloor n/d\rfloor$. Thus
\[
\frac nd = \left\lfloor\frac nd\right\rfloor + \frac{n\bmod d}{d}, \qquad 0\le\frac{n\bmod d}{d}<1.
\]
\end{definition}

\begin{example}
$6\bmod4=2$, \ $7\bmod4=3$, \ $6\bmod2=0$, \ $7\bmod2=1$, \ $6\bmod5=1$, \ $144\bmod100=44$, \ $174\bmod48=30$.
\end{example}

\begin{note}[Negative $n$]
The divisor $d$ is always positive, but $n$ may be negative; $q$ is still the largest integer with $qd\le n$. For example $-6\bmod5=4$: the largest multiple of $5$ that is $\le-6$ is $-10$ (so $q=-2$), and $-10+4=-6$. A programming language's \texttt{\%} operator need not agree with $\bmod$ on negative arguments.
\end{note}

\subsection{Parity}
\label{sec:7-4}
% CN 7.4

\begin{definition}[Parity]
The \emph{parity} of $x\in\Z$ is $x\bmod2$; $x$ is \emph{even} if this is $0$ and \emph{odd} if it is $1$.
\end{definition}

\begin{example}
The sum or difference of two odd numbers is even, while their product is odd.
\end{example}

\begin{note}[Parity as Boolean algebra]
Writing parities as $0,1$, the addition and multiplication tables mod $2$ are
\[
\begin{array}{c|cc}
+ & 0 & 1\\\hline
0 & 0 & 1\\
1 & 1 & 0
\end{array}
\qquad\qquad
\begin{array}{c|cc}
\times & 0 & 1\\\hline
0 & 0 & 0\\
1 & 0 & 1
\end{array}
\]
These are exactly the truth tables of $\oplus$ and $\wedge$ with $0=$ False and $1=$ True. So $\Z_2$ with $+,\times$ \emph{is} propositional logic with $\oplus,\wedge$, up to renaming. This is why logic gates can do arithmetic.
\end{note}

\subsection{The greatest common divisor}
\label{sec:7-5}
% CN 7.5

\begin{definition}[gcd, lcm]
For $a,b\in\Z$ not both $0$, $\gcd(a,b)$ is the greatest $d$ with $d\mid a$ and $d\mid b$. For $a,b\in\N$, $\lcmop(a,b)$ is the least positive common multiple of $a$ and $b$.
\end{definition}

\begin{example}
The positive divisors of $12$ are $1,2,3,4,6,12$ and those of $20$ are $1,2,4,5,10,20$. The common ones are $1,2,4$, so $\gcd(12,20)=4$.
\end{example}

\begin{note}[Why the gcd exists]
The set of common divisors contains $1$, and is bounded above because $d\mid n\ne0$ implies $d\le|n|$; so it has a greatest element. If $a=b=0$, every integer is a common divisor, which is why $\gcd(0,0)$ is undefined.
\end{note}

\begin{proposition}[Subtraction invariance]
\label{thm:gcdsub}
For $a,b\in\Z$ not both $0$, \ $\gcd(a,b)=\gcd(b,a-b)$.
\end{proposition}
\begin{proof}
We show that $(a,b)$ and $(b,a-b)$ have exactly the same common divisors; then their greatest common divisors are equal.

If $d\mid a$ and $d\mid b$, then $d\mid(a-b)$ by Proposition~\ref{thm:divlinear}, so $d$ is a common divisor of $b$ and $a-b$. Conversely, if $d\mid b$ and $d\mid(a-b)$, then $d\mid\bigl((a-b)+b\bigr)=a$, so $d$ is a common divisor of $a$ and $b$.
\end{proof}

\begin{note}[Repeated subtraction]
Applying Proposition~\ref{thm:gcdsub} repeatedly, always subtracting the smaller number from the larger, computes the gcd: $\gcd(20,12)=\gcd(12,8)=\gcd(8,4)=\gcd(4,4)=\gcd(4,0)=4$. The Euclidean algorithm (\S\ref{sec:7-6}) speeds this up.
\end{note}

\begin{proposition}[gcd--lcm identity (extra)]
\label{thm:gcdlcm}
For $a,b\in\N$, \ $\gcd(a,b)\cdot\lcmop(a,b)=ab$.
\end{proposition}
\begin{proof}
Write $a=\prod_i p_i^{e_i}$ and $b=\prod_i p_i^{f_i}$ over a common list of primes (Theorem~\ref{thm:fta}). Then $\gcd(a,b)=\prod_i p_i^{\min(e_i,f_i)}$ and $\lcmop(a,b)=\prod_i p_i^{\max(e_i,f_i)}$, and $\min(e_i,f_i)+\max(e_i,f_i)=e_i+f_i$ for each $i$.
\end{proof}

\subsection{The Euclidean algorithm}
\label{sec:7-6}
% CN 7.6

\begin{lemma}[Divisibility of integer combinations]\label{lem:div-comb}
Let $d,u,v \in \Z$. If $d \mid u$ and $d \mid v$, then $d \mid (su + tv)$ for all $s,t \in \Z$.
\end{lemma}
\begin{proof}
Write $u = dk_1$ and $v = dk_2$. Then $su + tv = d(sk_1 + tk_2)$, and $sk_1 + tk_2 \in \Z$.
\end{proof}

\begin{lemma}[Euclid's Lemma for the gcd]\label{lem:euclid-step}
Let $a \in \Z$ and $b \in \N$, and write $a = qb + r$ with $q = \lfloor a/b \rfloor$ and $0 \le r < b$. Then $\gcd(a,b) = \gcd(b,r)$.
\end{lemma}
\begin{proof}
As in Proposition~\ref{thm:gcdsub}, we show that $(a,b)$ and $(b,r)$ have the same common divisors.

If $d\mid a$ and $d\mid b$, then $d$ divides $r = 1\cdot a + (-q)\cdot b$ by Lemma~\ref{lem:div-comb}, so $d$ is a common divisor of $b$ and $r$. Conversely, if $d\mid b$ and $d\mid r$, then $d$ divides $a = q\cdot b + 1\cdot r$, so $d$ is a common divisor of $a$ and $b$.

The two sets of common divisors are equal, and nonempty and bounded above since $b>0$, so their greatest elements are equal: $\gcd(a,b)=\gcd(b,r)$.
\end{proof}

\begin{definition}[Euclidean Algorithm (EA)]
To find $\gcd(a,b)$:
\begin{enumerate}
    \item Input $a,b\in\N$; if $a<b$, swap them so that $a\ge b$.
    \item If $b\mid a$, output $b$.
    \item Otherwise output $\gcd(b,\,a\bmod b)$, computed by the same method.
\end{enumerate}
\end{definition}

\begin{note}[Division replaces repeated subtraction]
Subtracting $b$ from $a$ until the result drops below $b$ is exactly division with remainder. So one EA step does in one go what Proposition~\ref{thm:gcdsub} would do with many subtractions.
\end{note}

\begin{example}
$\gcd(1180,482)$: \ $1180=2\cdot482+216$, \ $482=2\cdot216+50$, \ $216=4\cdot50+16$, \ $50=3\cdot16+2$, \ $16=8\cdot2+0$. So $\gcd(1180,482)=2$.
\end{example}

\begin{example}[Worst case of the EA]
Consecutive Fibonacci numbers are the worst case: every quotient is $1$, so each remainder is only the previous Fibonacci number. For $(55,34)$:
\[
55=1\cdot34+21,\quad 34=1\cdot21+13,\quad 21=1\cdot13+8,\quad 13=1\cdot8+5,
\]
\[
8=1\cdot5+3,\quad 5=1\cdot3+2,\quad 3=1\cdot2+1,\quad 2=2\cdot1.
\]
So $\gcd(55,34)=1$. In the same way, consecutive Fibonacci numbers are always coprime (Exercise 7.12).
\end{example}

\subsection{The gcd and integer linear combinations}
\label{sec:7-7}
% CN 7.7

\begin{definition}[Integer linear combinations]
For $m,n\in\Z$, the set of all \emph{integer linear combinations} of $m$ and $n$ is
\[
m\Z+n\Z=\{xm+yn : x,y\in\Z\}.
\]
\end{definition}

\begin{example}
For $m=12$, $n=20$: \ $1\cdot12+0\cdot20=12$, \ $-1\cdot12+1\cdot20=8$, \ $2\cdot12-1\cdot20=4$, \ $-2\cdot12+1\cdot20=-4$, \ $-5\cdot12+3\cdot20=0$. In fact $12\Z+20\Z$ is exactly $4\Z$, the multiples of $\gcd(12,20)=4$.
\end{example}

\begin{lemma}
\label{thm:deltaZ}
Let $m,n\in\N$ and let $\Delta$ be the least positive element of $m\Z+n\Z$. Then $m\Z+n\Z=\Delta\Z$.
\end{lemma}
\begin{proof}
($\supseteq$) Write $\Delta=x_0m+y_0n$. For any $k\in\Z$, $k\Delta=(kx_0)m+(ky_0)n\in m\Z+n\Z$.

($\subseteq$) Let $z=xm+yn$ and divide: $z=q\Delta+r$ with $0\le r<\Delta$ (Theorem~\ref{thm:division}). Both $z$ and $q\Delta$ lie in $m\Z+n\Z$, which is closed under subtraction, so $r=z-q\Delta$ does too. If $r>0$ this would contradict the minimality of $\Delta$. So $r=0$ and $\Delta\mid z$.
\end{proof}

\begin{theorem}[Bézout]
\label{thm:bezout}
For $m,n\in\N$, $\gcd(m,n)$ is the least positive integer linear combination of $m$ and $n$. In particular, there exist $x,y\in\Z$ with $\gcd(m,n)=xm+yn$.
\end{theorem}
\begin{proof}
Let $\Delta$ be as in Lemma~\ref{thm:deltaZ}. Since $\Delta$ divides every element of $m\Z+n\Z$, in particular $m$ and $n$, it is a common divisor. Conversely, any common divisor $d$ of $m$ and $n$ divides every $xm+yn$ (Proposition~\ref{thm:divlinear}), in particular $\Delta$, so $d\le\Delta$. Hence $\Delta=\gcd(m,n)$.
\end{proof}

\begin{note}[Bézout coefficients are never unique]
Since $nm-mn=0$, adding $k$ times $(n,-m)$ to a solution gives another: $(x+kn)m+(y-km)n = xm+yn$. For example $4=2\cdot12-1\cdot20=(-3)\cdot12+2\cdot20$.
\end{note}

\subsection{The extended Euclidean algorithm}
\label{sec:7-8}
% CN 7.8

\begin{definition}[Extended Euclidean Algorithm (EEA)]
To find $\gcd(m,n)$ together with $x,y$ such that $\gcd(m,n)=xm+yn$:
\begin{enumerate}
    \item Input $m,n\in\N$; if $m<n$, swap them so that $m\ge n$.
    \item Initialise the triples $(a,x,y):=(m,1,0)$ and $(b,z,w):=(n,0,1)$.
    \item If $b=0$, output $(a,x,y)$.
    \item Let $q:=\lfloor a/b\rfloor$. Replace $(a,x,y)$ by $(b,z,w)$, and $(b,z,w)$ by $(a,x,y)-q\cdot(b,z,w)$ (using the old values). Go to step 3.
\end{enumerate}
\end{definition}

\begin{note}[The EA with extra bookkeeping]
The first entries of the triples evolve exactly as in the EA. The other two entries just record how each first entry is made up as an integer linear combination of $m$ and $n$.
\end{note}

\begin{theorem}[Correctness of the EEA]
\label{thm:eea}
At every stage the triples satisfy $a=xm+yn$ and $b=zm+wn$. On termination $a=\gcd(m,n)$, so the output gives Bézout coefficients.
\end{theorem}
\begin{proof}
Initially $m=1\cdot m+0\cdot n$ and $n=0\cdot m+1\cdot n$. If the invariant holds before an update, the new $a$-triple is the old $b$-triple, and the new $b$-triple satisfies $a-qb = (x-qz)m+(y-qw)n$, so the invariant still holds. Since the first entries follow the EA, at termination $a=\gcd(m,n)$.
\end{proof}

\begin{example}
Compute $\gcd(40,27)$ and express it as a combination. Each row $(t,x,y)$ satisfies $t=40x+27y$, and each new row is the row two above minus $q$ times the row above:
\[
\begin{array}{rrrl}
40 & 1 & 0 & \\
27 & 0 & 1 & \\
13 & 1 & -1 & (q=1)\\
1 & -2 & 3 & (q=2)\\
0 & 27 & -40 & (q=13)
\end{array}
\]
So $\gcd(40,27)=1=-2\cdot40+3\cdot27$. The row with $t=1$ already gives the answer; the last row is only a check.
\end{example}

\begin{example}[EEA with a nontrivial gcd]
For $\gcd(1180,482)$, with rows $t=1180x+482y$:
\[
\begin{array}{rrrl}
1180 & 1 & 0 & \\
482 & 0 & 1 & \\
216 & 1 & -2 & (q=2)\\
50 & -2 & 5 & (q=2)\\
16 & 9 & -22 & (q=4)\\
2 & -29 & 71 & (q=3)\\
0 & 241 & -590 & (q=8)
\end{array}
\]
So $\gcd(1180,482)=2=-29\cdot1180+71\cdot482$. Check: $71\cdot482-29\cdot1180=34222-34220=2$.
\end{example}

\begin{note}[Error check]
The signs in the last two columns alternate down the table. If they ever fail to alternate, an arithmetic mistake has been made.
\end{note}

\subsection{Coprimality}
\label{sec:7-9}
% CN 7.9

\begin{definition}[Coprime]
$a,b\in\Z$ are \emph{coprime} (\emph{relatively prime}) if $\gcd(a,b)=1$.
\end{definition}

\begin{note}[Coprime is not prime]
$21$ and $25$ are coprime, but neither is prime. Coprimality is easy to test (run the EA), whereas primality testing is a genuinely different and harder problem.
\end{note}

\begin{corollary}[Bézout characterisation of coprimality]
\label{thm:coprimebezout}
$m,n\in\Z$ are coprime if and only if $xm+yn=1$ for some $x,y\in\Z$.
\end{corollary}
\begin{proof}
By Theorem~\ref{thm:bezout}, $\gcd(m,n)$ is the least positive element of $m\Z+n\Z$. No positive integer is smaller than $1$, so this least element is $1$ if and only if $1\in m\Z+n\Z$.
\end{proof}

\begin{example}
$21$ and $25$ are coprime, and indeed $6\cdot21-5\cdot25=126-125=1$.
\end{example}

\begin{lemma}[Products preserve coprimality]\label{lem:coprimeprod}
If $\gcd(a,c)=\gcd(b,c)=1$, then $\gcd(ab,c)=1$.
\end{lemma}
\begin{proof}
By Corollary~\ref{thm:coprimebezout}, write $xa+uc=1$ and $yb+vc=1$. Multiplying these gives $1=(xy)\,ab+(xav+uyb+uvc)\,c$, so $\gcd(ab,c)=1$ by Corollary~\ref{thm:coprimebezout} again.
\end{proof}

\bigskip
\begin{theorem}[Euclid's lemma]
\label{thm:euclidlemma}
Let $p$ be prime and $a,b\in\Z$. If $p\mid ab$, then $p\mid a$ or $p\mid b$.
\end{theorem}
\begin{proof}
Assume $p\mid ab$ and $p\nmid a$; we show $p\mid b$. Since $p$ is prime and $p\nmid a$, $\gcd(p,a)=1$, so by Corollary~\ref{thm:coprimebezout} there are $x,y\in\Z$ with $xp+ya=1$. Multiplying by $b$ gives $xpb+yab=b$. The first term is a multiple of $p$, and the second is a multiple of $ab$, hence of $p$. So $p\mid b$.
\end{proof}

\begin{note}[Primality is needed]
$3\mid(8\times9)$ forces $3\mid8$ or $3\mid9$ (here $3\mid9$). But for non-prime divisors this fails: $6\mid(3\times4)$, yet $6\nmid3$ and $6\nmid4$.
\end{note}

\begin{proposition}[Coprime divisor rule]
\label{thm:coprimedivides}
If $a\mid n$, $b\mid n$ and $\gcd(a,b)=1$, then $ab\mid n$.
\end{proposition}
\begin{proof}
Write $xa+yb=1$ (Corollary~\ref{thm:coprimebezout}) and $n=a\alpha=b\beta$. Then
\[
n = n(xa+yb) = xa(b\beta)+yb(a\alpha) = ab(x\beta+y\alpha). \qedhere
\]
\end{proof}

\begin{theorem}[Fundamental Theorem of Arithmetic]
\label{thm:fta}
Every positive integer has a factorisation into primes that is unique up to order.
\end{theorem}
\begin{proof}
Existence is Theorem~\ref{thm:primefactexists}. For uniqueness, suppose, for contradiction, that some integer has two different prime factorisations, and let $n$ be the smallest such. Over the primes $p_1,\dots,p_k$ appearing in either factorisation, write $n=p_1^{e_1}\cdots p_k^{e_k}=p_1^{f_1}\cdots p_k^{f_k}$, with $e_i\ne f_i$ for some $i$.

If some prime $p_j$ appears in both factorisations, divide both by $p_j^{\min(e_j,f_j)}$. The results are still two different factorisations, of the smaller number $n/p_j^{\min(e_j,f_j)}$, contradicting the minimality of $n$.

So the two factorisations share no prime. Let $p$ occur in the first. Then $p\mid n$, while $n$ is also a product of primes $q_1,\dots,q_l$, none equal to $p$. Applying Theorem~\ref{thm:euclidlemma} repeatedly gives $p\mid q_t$ for some $t$, so $p=q_t$, a contradiction.
\end{proof}

\subsection{Modular arithmetic}
\label{sec:7-10}
% CN 7.10

\begin{note}[Where remainders arise]
Often only remainders matter. Days of the week are remainders mod $7$ of the number of days since a reference date; hours in 24-hour time are remainders mod $24$; the state of a toggle switch ($0$ off, $1$ on) is the number of presses mod $2$.
\end{note}

\begin{definition}[Congruence]
For $n\in\N$ and $a,b\in\Z$, $a$ is \emph{congruent to $b$ modulo $n$}, written $a\equiv b\pmod n$, if $n\mid(a-b)$.
\end{definition}

\begin{proposition}
\label{thm:congremainder}
$a\equiv b\pmod n$ if and only if $a\bmod n=b\bmod n$.
\end{proposition}
\begin{proof}
Write $a=q_an+r_a$ and $b=q_bn+r_b$ with $0\le r_a,r_b\le n-1$. Then $a-b=(q_a-q_b)n+(r_a-r_b)$, so $n\mid(a-b)$ if and only if $n\mid(r_a-r_b)$. Since $|r_a-r_b|\le n-1<n$, this holds if and only if $r_a=r_b$.
\end{proof}

\begin{proposition}
\label{thm:congequiv}
For every $n\in\N$, congruence modulo $n$ is an equivalence relation on $\Z$.
\end{proposition}
\begin{proof}
\textbf{Reflexive}: $a-a=0=0\cdot n$. \ \textbf{Symmetric}: if $a-b=kn$ then $b-a=(-k)n$. \ \textbf{Transitive}: if $a-b=kn$ and $b-c=ln$ then $a-c=(k+l)n$.
\end{proof}

\begin{definition}[Residue classes and $\Z_n$]
The equivalence classes of congruence mod $n$ are the \emph{residue classes} $[a]=a+n\Z=\{a+kn : k\in\Z\}$. There are exactly $n$ of them: $[0],[1],\dots,[n-1]$. We write $\Z_n$ for $\{0,1,\dots,n-1\}$ with $+,-,\times$ followed by reduction mod $n$.
\end{definition}

\begin{note}[Equation in $\Z_n$ vs.\ congruence]
In $\Z_4$, $2+3=1$, since $5\bmod4=1$. This is different from the congruence $2+3\equiv1\pmod4$: the congruence only says that $2+3$ and $1$ lie in the same class, so $2+3\equiv37\pmod4$ is also true, while ``$2+3=37$ in $\Z_4$'' makes no sense ($37\notin\Z_4$).
\end{note}

\begin{theorem}[Congruences respect arithmetic]
\label{thm:congarith}
If $a\equiv b\pmod n$ and $c\equiv d\pmod n$, then
\[
a+c\equiv b+d, \qquad a-c\equiv b-d, \qquad ac\equiv bd \pmod n.
\]
\end{theorem}
\begin{proof}
Write $a-b=kn$ and $c-d=ln$. Then $(a+c)-(b+d)=(k+l)n$ and $(a-c)-(b-d)=(k-l)n$. For the product, $ac-bd = c(a-b)+b(c-d) = (ck+bl)n$.
\end{proof}

\begin{note}[Reduce as you go]
By Theorem~\ref{thm:congarith}, in any calculation using only $+,-,\times$, you can reduce mod $n$ at \emph{every} step instead of only at the end. Division and exponentiation need separate treatment (\S\ref{sec:7-11}, \S\ref{sec:7-14}).
\end{note}

\begin{example}[Checking a product mod 9]
Is $84\times32=2678$? Since $10\equiv1\pmod9$, a number is congruent mod $9$ to its digit sum. So $84\equiv12\equiv3$ and $32\equiv5$, giving $84\times32\equiv15\equiv6$; but $2678\equiv23\equiv5$. The two sides differ mod $9$, so the calculation is wrong (in fact $84\times32=2688$). The check is one-sided: agreement mod $9$ would not prove the answer correct.
\end{example}

\subsection{Modular inverses}
\label{sec:7-11}
% CN 7.11

\begin{definition}[Modular inverse, $\Z_n^*$]
$x^{-1}\in\Z_n$ is a \emph{multiplicative inverse} of $x\in\Z_n$ if $xx^{-1}\equiv1\pmod n$. We write $\Z_n^*$ for the set of elements of $\Z_n$ that have an inverse.
\end{definition}

\begin{example}
In $\Z_7$: $2^{-1}=4$ (since $2\cdot4=8\equiv1$), and $1^{-1}=1$, $3^{-1}=5$, $4^{-1}=2$, $5^{-1}=3$, $6^{-1}=6$. Every nonzero element has an inverse, so $\Z_7^*=\{1,2,3,4,5,6\}$.

In $\Z_6$, $2$ has no inverse: $2z$ is always even, so it is never $1,7,13,\dots$. Similarly $3$ and $4$ have none, while $5\cdot5=25\equiv1$. So $\Z_6^*=\{1,5\}$.
\end{example}

\begin{theorem}[Existence of inverses]
\label{thm:inverse}
$x\in\Z_n$ has an inverse in $\Z_n$ if and only if $\gcd(x,n)=1$.
\end{theorem}
\begin{proof}
($\Rightarrow$) If $xx^{-1}\equiv1\pmod n$, then $xx^{-1}+kn=1$ for some $k\in\Z$, so $x$ and $n$ are coprime by Corollary~\ref{thm:coprimebezout}.

($\Leftarrow$) If $\gcd(x,n)=1$, Corollary~\ref{thm:coprimebezout} gives $y,z\in\Z$ with $yx+zn=1$. Let $y'=y\bmod n$, so $y=ln+y'$. Then $y'x+(lx+z)n=1$, so $y'x\equiv1\pmod n$ with $y'\in\Z_n$.
\end{proof}

\begin{note}[Finding inverses with the EEA]
The proof is constructive: run the EEA on $(x,n)$ to get $yx+zn=1$; then $x^{-1}=y\bmod n$. This is the main practical use of the EEA.
\end{note}

\begin{example}[Inverse computation]
Find $17^{-1}$ in $\Z_{101}$. Rows $(t,x,y)$ satisfy $t=101x+17y$:
\[
\begin{array}{rrrl}
101 & 1 & 0 &\\
17 & 0 & 1 &\\
16 & 1 & -5 & (q=5)\\
1 & -1 & 6 & (q=1)
\end{array}
\]
So $1=-1\cdot101+6\cdot17$, giving $17^{-1}=6$ in $\Z_{101}$. Check: $6\cdot17=102\equiv1\pmod{101}$.
\end{example}

\subsection{The Euler totient function}
\label{sec:7-12}
% CN 7.12

\begin{definition}[Euler totient function]\label{def:totient}
For $n\in\N$, $\varphi(n)$ is the number of integers $x$ with $1\le x\le n$ and $\gcd(x,n)=1$. By Theorem~\ref{thm:inverse}, $\varphi(n)=|\Z_n^*|$.
\end{definition}

\begin{example}
$\varphi(6)=|\{1,5\}|=2$ and $\varphi(7)=|\{1,\dots,6\}|=6$.
\end{example}

\begin{note}[Range of $x$]
For $n\ge2$ this agrees with CN's range $0<x<n$, since $\gcd(n,n)=n>1$. Including $x=n$ gives $\varphi(1)=1$, which keeps multiplicativity valid when one factor is $1$.
\end{note}

\begin{proposition}[Totient of a prime power]
\label{thm:totientpp}
For $p$ prime and $m\in\N$, \ $\varphi(p^m)=p^{m-1}(p-1)$. In particular $\varphi(p)=p-1$.
\end{proposition}
\begin{proof}
An integer $x$ with $1\le x\le p^m$ fails to be coprime to $p^m$ exactly when $p\mid x$. The multiples of $p$ in this range are $p,2p,\dots,p^{m-1}\cdot p$, so there are $p^{m-1}$ of them. Hence $\varphi(p^m)=p^m-p^{m-1}=p^{m-1}(p-1)$.
\end{proof}

\begin{theorem}[Multiplicativity]
\label{thm:totientmult}
If $\gcd(a,b)=1$ then $\varphi(ab)=\varphi(a)\varphi(b)$.
\end{theorem}

\begin{theorem}[General formula]
\label{thm:totientformula}
If $n=p_1^{m_1}\cdots p_k^{m_k}$ with distinct primes $p_i$, then
\[
\varphi(n)=\prod_{i=1}^{k}p_i^{m_i-1}(p_i-1) = n\prod_{i=1}^{k}\left(1-\frac1{p_i}\right).
\]
In particular, for distinct primes $p,q$, \ $\varphi(pq)=(p-1)(q-1)$.
\end{theorem}
\begin{proof}
Powers of distinct primes are coprime, so Theorem~\ref{thm:totientmult} splits $\varphi(n)$ into a product of the $\varphi(p_i^{m_i})$, each given by Proposition~\ref{thm:totientpp}. The second form follows from $p_i^{m_i-1}(p_i-1)=p_i^{m_i}(1-1/p_i)$.
\end{proof}

\begin{example}
$\varphi(12)=\varphi(4)\varphi(3)=2\cdot2=4$; indeed the integers in $[1,12]$ coprime to $12$ are $1,5,7,11$. And $\varphi(25)=5\cdot4=20$.
\end{example}

\begin{note}[Needs the factorisation]
The formula needs the prime factorisation of $n$, and no essentially easier way to compute $\varphi(n)$ is known. The RSA cryptosystem relies on this.
\end{note}

\subsection{Fast exponentiation}
\label{sec:7-13}
% CN 7.13

\begin{definition}[Binary (square-and-multiply) exponentiation]
Write $m=\sum_{i}b_i2^i$ in binary. Then $a^m=\prod_{i \,:\, b_i=1}a^{2^i}$, where the powers $a^{2^i}$ are obtained by repeated squaring: $a^{2^{i+1}}=\bigl(a^{2^i}\bigr)^2$.
\end{definition}

\begin{proposition}[Cost]
\label{thm:fastexpcost}
For $m\ge1$, binary exponentiation computes $a^m$ using at most $2\lfloor\log_2 m\rfloor$ multiplications.
\end{proposition}
\begin{proof}
Producing $a^2,a^4,\dots,a^{2^{\lfloor\log_2m\rfloor}}$ takes $\lfloor\log_2m\rfloor$ squarings. The binary expansion of $m$ has $\lfloor\log_2m\rfloor+1$ digits, so at most that many ones, and combining the corresponding powers takes at most $\lfloor\log_2m\rfloor$ further multiplications.
\end{proof}

\begin{example}[Square-and-multiply]
Compute $3^{36}$. Since $36=100100_2=2^5+2^2$, we have $3^{36}=3^{32}\cdot3^4$. Squaring repeatedly gives $3^2=9$, $3^4=81$, $3^8=6561$, $3^{16}=43046721$, $3^{32}=1853020188851841$, and one final multiplication by $3^4=81$ gives $3^{36}$. That is six multiplications, versus $35$ for the naive method.
\end{example}

\subsection{Modular exponentiation}
\label{sec:7-14}
% CN 7.14

\begin{note}[Two savings]
Modular exponentiation is exponentiation in $\Z_n$. Two separate savings apply: reduce each intermediate product mod $n$ (Theorem~\ref{thm:congarith}), and, when $\gcd(a,n)=1$, reduce the \emph{exponent} mod $\varphi(n)$ (Corollary~\ref{cor:expred} below).
\end{note}

\begin{theorem}[Euler]
\label{thm:euler}
If $\gcd(x,n)=1$, then $x^{\varphi(n)}\equiv1\pmod n$.
\end{theorem}
\begin{proof}
Let $\Z_n^*=\{y_1,\dots,y_{\varphi(n)}\}$ and consider $xy_1,\dots,xy_{\varphi(n)}$. These lie in $\Z_n^*$ by Lemma~\ref{lem:coprimeprod}. They are also distinct mod $n$: if $xy_i\equiv xy_j$, multiplying by $x^{-1}$ (which exists by Theorem~\ref{thm:inverse}) gives $y_i=y_j$. So multiplying by $x$ just permutes $\Z_n^*$, and the two lists have the same product:
\[
y_1y_2\cdots y_{\varphi(n)} \equiv (xy_1)(xy_2)\cdots(xy_{\varphi(n)}) = x^{\varphi(n)}\,y_1y_2\cdots y_{\varphi(n)} \pmod n.
\]
The product $y_1\cdots y_{\varphi(n)}$ lies in $\Z_n^*$, so it has an inverse; multiplying both sides by it gives $x^{\varphi(n)}\equiv1\pmod n$.
\end{proof}

\begin{corollary}[Fermat's Little Theorem]
\label{thm:flt}
If $p$ is prime and $p\nmid x$, then $x^{p-1}\equiv1\pmod p$.
\end{corollary}
\begin{proof}
Take $n=p$ in Theorem~\ref{thm:euler}, and use $\varphi(p)=p-1$ (Proposition~\ref{thm:totientpp}).
\end{proof}

\begin{corollary}[Exponent reduction]\label{cor:expred}
If $\gcd(a,n)=1$ and $x\equiv y\pmod{\varphi(n)}$ with $x,y\ge0$, then $a^x\equiv a^y\pmod n$.
\end{corollary}
\begin{proof}
Taking $x\ge y$, write $x=y+k\varphi(n)$. Then $a^x=a^y\bigl(a^{\varphi(n)}\bigr)^k\equiv a^y\cdot1^k=a^y\pmod n$ by Theorem~\ref{thm:euler}.
\end{proof}

\begin{note}[The coprimality hypothesis is essential]
For $n=4$: $\varphi(4)=2$ and $3\equiv1\pmod2$, yet $2^3=8\equiv0$ while $2^1=2\pmod4$.
\end{note}

\begin{example}
Simplify $2^9\bmod7$. Since $\varphi(7)=6$, the exponent reduces mod $6$: $9\bmod6=3$. So $2^9\equiv2^3=8\equiv1\pmod7$.
\end{example}

\begin{example}[Exponent reduction]
Compute $3^{36}\bmod25$. Here $\varphi(25)=20$ and $\gcd(3,25)=1$, so $3^{36}\equiv3^{36\bmod20}=3^{16}\pmod{25}$. Squaring repeatedly mod $25$: $3^2=9$, $3^4=81\equiv6$, $3^8\equiv6^2=36\equiv11$, $3^{16}\equiv11^2=121\equiv21$. So $3^{36}\equiv21\pmod{25}$, using four multiplications and no number larger than $121$.
\end{example}

\subsection{Primitive roots}
\label{sec:7-15}
% CN 7.15

\begin{definition}[Primitive root]
Let $\gcd(x,n)=1$. Then $x$ is a \emph{primitive root} of $n$ if $x^k\not\equiv1\pmod n$ for all $1\le k<\varphi(n)$, i.e.\ $\varphi(n)$ is the least positive exponent with $x^{\varphi(n)}\equiv1$. Equivalently, the powers $x,x^2,x^3,\dots$ run through all of $\Z_n^*$; we say $x$ \emph{generates} $\Z_n^*$.
\end{definition}

\begin{example}[Primitive root and a non-example]
For $n=7$, $\varphi(7)=6$ and $\Z_7^*=\{1,\dots,6\}$. The powers of $3$ are
\[
\begin{array}{c|cccccc}
k & 1&2&3&4&5&6\\\hline
3^k\bmod7 & 3&2&6&4&5&1
\end{array}
\]
All six elements appear, and $1$ is first reached at $k=6$, so $3$ is a primitive root of $7$. By contrast, the powers of $2$ are $2,4,1,\dots$, reaching $1$ at $k=3<6$; so $2$ generates only $\{1,2,4\}$ and is not a primitive root.
\end{example}

\begin{note}[Primitive roots have the widest range]
Among the members of $\Z_n^*$, the primitive roots are exactly those whose powers take the greatest possible number of values. This is what makes them useful in cryptography (\S\ref{sec:7-17}).
\end{note}

\begin{theorem}[Which moduli have primitive roots]
\label{thm:primrootexists}
$n$ has a primitive root if and only if $n\in\{1,2,4\}$, or $n=p^k$ or $n=2p^k$ for an odd prime $p$ and $k\in\N$.
\end{theorem}

\begin{example}[No primitive root]
$8$ has no primitive root. Here $\varphi(8)=4$ and $\Z_8^*=\{1,3,5,7\}$, but $3^2=9$, $5^2=25$ and $7^2=49$ are all $\equiv1\pmod8$, so every element reaches $1$ by exponent $2<4$. Elements outside $\Z_8^*$ cannot help: the powers of $2$ are $2,4,0,0,\dots$ and never reach $1$. This agrees with Theorem~\ref{thm:primrootexists}, since $8=2^3$ is a power of the \emph{even} prime.
\end{example}

\bigskip
\begin{lemma}[Order lemma]
\label{lem:order}
Let $\gcd(x,n)=1$ and let $t\ge1$ be the least positive integer with $x^t\equiv1\pmod n$ (it exists by Theorem~\ref{thm:euler}). Then for every $m\ge0$, \ $x^m\equiv1\pmod n$ if and only if $t\mid m$.
\end{lemma}
\begin{proof}
($\Leftarrow$) If $m=qt$ then $x^m=(x^t)^q\equiv1$.

($\Rightarrow$) Divide: $m=qt+r$ with $0\le r<t$. Then $1\equiv x^m=(x^t)^q\,x^r\equiv x^r\pmod n$. By the minimality of $t$, $r$ cannot be positive, so $r=0$ and $t\mid m$.
\end{proof}

\begin{theorem}[Number of primitive roots]
\label{thm:primrootcount}
If $n$ has a primitive root, then it has exactly $\varphi(\varphi(n))$ of them.
\end{theorem}
\begin{proof}
Let $a$ be a primitive root of $n$ and write $\phi=\varphi(n)$.

\textbf{Every element of $\Z_n^*$ is $a^k$ for a unique $k\in\{1,\dots,\phi\}$.} Each $a^k$ lies in $\Z_n^*$ (Lemma~\ref{lem:coprimeprod}). If $a^k\equiv a^l$ with $1\le k<l\le\phi$, multiplying by $(a^{-1})^k$ gives $a^{l-k}\equiv1$ with $1\le l-k<\phi$, contradicting the definition of a primitive root. So $a^1,\dots,a^\phi$ are $\phi=|\Z_n^*|$ distinct elements of $\Z_n^*$, i.e.\ all of it.

\textbf{If $d=\gcd(k,\phi)>1$, then $a^k$ is not a primitive root}, since $(a^k)^{\phi/d}=(a^{\phi})^{k/d}\equiv1$ with $1\le\phi/d<\phi$.

\textbf{If $\gcd(k,\phi)=1$, then $a^k$ is a primitive root.} Suppose $(a^k)^j=a^{kj}\equiv1$ with $j\ge1$. Since $\phi$ is the least positive exponent sending $a$ to $1$, Lemma~\ref{lem:order} gives $\phi\mid kj$. Write $x\phi+yk=1$ (Corollary~\ref{thm:coprimebezout}); then $j=x\phi j+y\,kj$ is a sum of multiples of $\phi$, so $\phi\mid j$ and $j\ge\phi$.

So the primitive roots of $n$ correspond one-to-one to the $k\in\{1,\dots,\phi\}$ with $\gcd(k,\phi)=1$, and there are $\varphi(\phi)=\varphi(\varphi(n))$ of these.
\end{proof}

\begin{example}
For $n=7$: $\varphi(\varphi(7))=\varphi(6)=2$. Indeed the primitive roots are $3^1=3$ and $3^5\equiv5$, while $3^2\equiv2$, $3^3\equiv6$, $3^4\equiv4$ are not, as their exponents $2,3,4$ share a factor with $6$.
\end{example}

\begin{note}[Finding primitive roots is hard]
No fast algorithm is known for finding a primitive root of a given $n$.
\end{note}

\subsection{One-way functions}
\label{sec:7-16}
% CN 7.16

\begin{definition}[One-way function, informal]
A function $f$ is \emph{one-way} if it is easy to compute but hard to invert, for almost all inputs.
\end{definition}

\begin{note}[None are proved]
No function has been \emph{proved} to be one-way; doing so would resolve $P\ne NP$ (\S\ref{sec:6-9}). Several functions are believed to be one-way and are used as such.
\end{note}

\begin{note}[Application: password storage]
Instead of storing each password $P_i$, a system stores $(U_i,f(P_i))$ for a one-way $f$. A login with password $P$ is accepted iff $f(P)=f(P_i)$, which is fast to check. An attacker who steals $f(P_i)$ still faces the inversion problem. This is not encryption: no key is involved, and nobody is meant to recover $P_i$.
\end{note}

\subsection{Modular exponentiation with fixed base}
\label{sec:7-17}
% CN 7.17

\begin{definition}[Modular exponentiation with fixed base]
Fix $n\in\N$ and a base $a<n$. The function is $x\longmapsto a^x\bmod n$, for $x<\varphi(n)$.
\end{definition}

\begin{definition}[Discrete Logarithm problem]
Fix $n\in\N$ and a primitive root $a$ of $n$. Given $y\in\Z_n^*$, find $x\le\varphi(n)$ with $a^x\equiv y\pmod n$. We write $x=\log_a y$.
\end{definition}

\begin{example}
With $n=7$ and $a=3$, the table in \S\ref{sec:7-15} gives $\log_3 6=3$ and $\log_3 4=4$.
\end{example}

\begin{note}[Choosing the parameters]
The base is chosen to be a primitive root because then $a^x$ takes the widest range of values, which maximises the search space for an attacker. Discrete Logarithm is believed to have no fast algorithm, and in practice is about as hard as integer factorisation. Since $\varphi(n)$ is largest when $n$ is prime, the best candidate one-way function uses a large prime $p$ and a primitive root $a$ of $p$.
\end{note}

\subsection{Diffie--Hellman key agreement}
\label{sec:7-18}
% CN 7.18

\begin{definition}[Diffie--Hellman scheme]
The public global parameters are a large prime $p$ and a primitive root $a$ of $p$. Each user picks a private random number $x\in\{1,\dots,p-1\}$ and publishes $y=a^x\bmod p$. Users $A$ and $B$, with private numbers $x_A,x_B$ and public numbers $y_A,y_B$, compute
\[
k_{AB} := y_B^{\,x_A}\bmod p, \qquad k_{BA} := y_A^{\,x_B}\bmod p.
\]
\end{definition}

\begin{proposition}[Agreement]
\label{thm:dhagree}
$k_{AB}=k_{BA}$.
\end{proposition}
\begin{proof}
In $\Z_p^*$, \ $k_{AB}=y_B^{\,x_A}=(a^{x_B})^{x_A}=a^{x_Ax_B}=(a^{x_A})^{x_B}=y_A^{\,x_B}=k_{BA}$.
\end{proof}

\begin{example}[Diffie--Hellman with small numbers]
Take $p=11$ and $a=2$, a primitive root of $11$ (Exercise 7.17). Alice picks $x_A=3$ and Bob picks $x_B=6$, so $y_A=2^3\bmod11=8$ and $y_B=2^6\bmod11=64\bmod11=9$. They exchange $y_A$ and $y_B$, and compute
\[
k_{AB}=9^3\bmod11=729\bmod11=3, \qquad k_{BA}=8^6\bmod11=262144\bmod11=3.
\]
Both equal $a^{x_Ax_B}=2^{18}\equiv3\pmod{11}$, which an eavesdropper who sees only $8$ and $9$ cannot easily compute.
\end{example}

\begin{note}[Agreement, not encryption]
Diffie--Hellman is a key \emph{agreement} scheme, not a cryptosystem: neither party can use it to send a message of their choosing. The shared $k_{AB}$ is then used as the key of a conventional secret-key cryptosystem. Its value is that no secure channel is needed to set up the shared key, and neither party alone decides it.
\end{note}

\bigskip
\begin{definition}[Diffie--Hellman problem]
Given $p$, $a$, $a^{x_A}\bmod p$ and $a^{x_B}\bmod p$, find $a^{x_Ax_B}\bmod p$.
\end{definition}

\begin{proposition}
\label{thm:dhreduce}
Any algorithm for Discrete Logarithm yields an algorithm for the Diffie--Hellman problem, with only polynomial extra work.
\end{proposition}
\begin{proof}
Give $p,a,a^{x_A}$ to the Discrete Logarithm algorithm to recover $x_A$. Then compute $(a^{x_B})^{x_A}=a^{x_Ax_B}\bmod p$ by fast modular exponentiation (\S\ref{sec:7-14}).
\end{proof}

\begin{note}[Open problem]
Whether the reduction also goes the other way (an efficient algorithm for the Diffie--Hellman problem giving one for Discrete Logarithm) is widely believed but unproven.
\end{note}
